
\chapter{Introducción}
\label{cap:introduccion}

Este capítulo presenta el contexto del trabajo, la motivación y los objetivos, y describe cómo se organiza el resto del documento. Los capítulos de esta plantilla incluyen ejemplos de uso: consulta su código fuente para ver cómo se hace cada cosa y sustituye el texto por el tuyo.

\section{Motivación}

Explica aquí el problema que se aborda y por qué es relevante. Los acrónimos se definen en \texttt{anexos/acronimos.tex} y se usan con \verb|\gls{etiqueta}|: la primera vez se escribe la forma completa, como en \gls{tfg} o \gls{api}, y las siguientes solo la sigla: \gls{tfg}, \gls{api}. La lista de acrónimos al final del documento solo incluye los que se han usado.

\section{Objetivos}

Las listas se escriben con \texttt{itemize} o \texttt{enumerate}:

\begin{enumerate}
  \item Objetivo principal del trabajo.
  \item Primer objetivo específico.
  \item Segundo objetivo específico.
\end{enumerate}

\section{Estructura del documento}

Para referirse a otras partes del documento se usa \verb|\cref{etiqueta}|, que añade automáticamente el tipo de elemento: el \cref{cap:estado-arte} revisa el estado del arte, el \cref{cap:analisis} describe los objetivos y la metodología, el \cref{cap:diseno} recoge el diseño y el \cref{cap:conclusiones} presenta las conclusiones. El \cref{anx:complementario} contiene material complementario.
